\scanpage{120}
% archive sheet number: 119

``compatible avec $\otimes$ èommutatif'' est un foncteur et un $c$, compatible avec la donnée des$\unsure{\otimes|\deleted{x}}$ $\handwritten{\unsure{\mu|\psi}}$ .

Quand on a une loi $\otimes$ avec donnée a) et b), il y a lieu de postuler une compatibilité entre les deux (dans le cadre del'axiomatique 1), est-ce une conséquence des autres axiomes ? je \deleted{crains}\inserted{\handwritten{pense}} que non, et qu'il faille le rajouter). Quand on a une loi $\otimes$ avec donnée a) et \unsure{v|c}), il y a lieu de postuler une compatibilité entre les deux. De même pour une loi $\otimes$ avec b) et c). Quand on a à la fois a)b)c), le paquet \deleted{\unsure{d'ax}} de toutes les conditions de compatibilité sont celles énoncées plus haut \inserted{dans 1º} a) à d) . De même, un Foncteur``compatible avec $\otimes$ associatif à \unsure{unité|mité} bilatère'' est un foncteur $F$, $+$ un $c$, compatible séparément avec associativité et unité bilatèrité, et de même pour tout autre paquet de données supplémentaires parmi a),b),c). Enfin, la notion d'homomorphisme de \deleted{\unsure{foncteu}} $\otimes$-foncteurs ne dépend pas des structures supplémentaires a)\deleted{b)} ou c) qu'on peut envisager\inserted{\handwritten{(mais bien de b)}} , s'exprimant exclusivement à l'aide de la donnée de $c$\deleted{\unsure{)}} \deleted{(sous réserve de vérification indiquée dans 3º)}. - Quan\deleted{d}\handwritten{t} a la \deleted{structure bilinéaire, il y a lieu de}\inserted{\deleted{\unsure{biaddiv|biadditiv}}} condition d'additivité sur les catégories et de biadditivité sur $\otimes$, il n'y a lieu de l'introduire qu'en dernier lieu.

\margin{Attention si l'associativité ne doit pas rappliquer ici aussi}
\deleted{5)}\inserted{\handwritten{5)}} NB Si on a $\otimes$ avec unité bilatère, alors pour deux endomorphismes $u,v$ de $\underline{1}$, on a $u\otimes v\overset{\handwritten{!}}{=}uv=vu$. Donc sous la condition de commutativité\deleted{\unsure{x|,}} c), $\unsure{\mathbb{E}|\mathrm{E}}\mathrm{nd}(\underline{1})$ est un monoïde commutatif. \inserted{\deleted{il opère sur tous}}\deleted{Si $\underline{C}$ est additif, c'est} Sans condition de comm. ce monoïde opère à gauche et à droite sur tout objet de $\underline{C}$; les deux opérations coïncident moyennant commutativité. Si \deleted{la structure}\inserted{$C$} est add\unsure{i|o}tive, $\mathrm{End}(\underline{1})$ est un anneau, qui opère à gauche et à droite sur chaque objet de $C$. Dans le cas commutatif, $\mathrm{End}(\underline{1})=\unsure{\mathbb{k}|k}$ est donc un anneau commytatif et la catégorie $C$ est $k$-linéaire de façon naturelle\handwritten{, $\otimes$ est $k$-bilinéaire.}
